\documentclass[a4paper]{article}

\usepackage[spanish]{babel}
\usepackage{graphicx}
\usepackage{amsmath, amssymb, mathrsfs}
\usepackage[margin=2cm]{geometry}
\usepackage{fancyhdr}
\usepackage{enumerate}
\usepackage[shortlabels]{enumitem}
\usepackage{parskip}
\usepackage[most]{tcolorbox}
\usepackage[hidelinks]{hyperref}
\usepackage{float}
\usepackage{booktabs}



% cabecera
\pagestyle{fancy}
\fancyhead[l]{Física Matemática I}
\fancyhead[c]{ParcialOtroCurso}

\renewcommand{\headrulewidth}{0.2pt} % linea horizontal

\begin{document}

\begin{titlepage}
  \centering
  {\Large Universidad de Antioquia\par}
  \vspace{2cm}
  {\huge\bfseries Física Matemática I\par}
  \vspace{0.4cm}
  {\Large ParcialOtroCurso \par}
  \vspace{2.5cm}
  {\large Autor:\par}
  \vspace{0.2cm}
  {\large Daniel Soto Villada\par}
  {\large Estudiante de Astronomía\par}
  \vfill
  {\large \today\par}
\end{titlepage}

\newpage
\tableofcontents
\newpage

\section{Problema 1}
Considere una esfera sólida a una temperatura en equilibrio. La distribución de calor dentro de la esfera de radio $\alpha$ (independiente del tiempo) cuando no hay fuentes presentes está dada por el valor de problema inicial

$$\nabla^2 u(r, \theta, \phi) = 0, \quad 0 \leq r \leq \alpha, \quad 0 \leq \phi \leq 2\pi, \quad 0 \leq \theta \leq \pi,$$

y la temperatura en la superficie está dada por

$$u(\alpha, \theta, \phi) = f(\theta, \phi) = \begin{cases} 2, & 0 < \phi < 2\pi, \quad 0 < \theta < \frac{\pi}{2}, \\ -1, & 0 < \phi < 2\pi, \quad \frac{\pi}{2} < \theta < \pi, \end{cases}$$

proponga una solución de la forma

$$u(\alpha, \theta, \phi) = R(r)\Theta(\theta)\Phi(\phi) \approx R(r)Y_{\ell, m}(\theta, \phi),$$

y halle una solución sujeta a éstas condiciones para el interior de la esfera.

\subsection{Solución}

Para resolver la ecuación de Laplace $\nabla^2 u = 0$ dentro de una esfera de radio $\alpha$ con simetría azimutal (ya que la condición de contorno $f(\theta, \phi)$ es independiente de $\phi$), la solución general se expresa mediante armónicos esféricos simplificados a polinomios de Legendre $P_\ell(\cos\theta)$:

$$u(r, \theta) = \sum_{\ell=0}^{\infty} A_\ell r^\ell P_\ell(\cos\theta)$$

Determinación de los coeficientes $A_\ell$

En la superficie $r = \alpha$, tenemos $u(\alpha, \theta) = f(\theta)$. Aplicamos la ortogonalidad de los polinomios de Legendre:

$$A_\ell \alpha^\ell = \frac{2\ell+1}{2} \int_{0}^{\pi} f(\theta) P_\ell(\cos\theta) \sin\theta \, d\theta$$

Realizando el cambio de variable $x = \cos\theta$ ($dx = -\sin\theta \, d\theta$), los límites pasan de $[0, \pi]$ a $[1, -1]$:

$$A_\ell \alpha^\ell = \frac{2\ell+1}{2} \left[ \int_{0}^{1} (2) P_\ell(x) \, dx + \int_{-1}^{0} (-1) P_\ell(x) \, dx \right]$$

Cálculo de las integrales

Los valores de estas integrales dependen de la paridad de $\ell$:

* Para $\ell = 0$: $A_0 = \frac{1}{2} [2(1) - 1(1)] = \frac{1}{2}$.

* Para $\ell$ impar: Los términos se suman, resultando en $A_\ell \alpha^\ell = \frac{2\ell+1}{2} \cdot 3 \int_{0}^{1} P_\ell(x) \, dx$.

* Para $\ell$ par ($\ell > 0$): La integral es cero.

Usando la propiedad $\int_{0}^{1} P_\ell(x) \, dx = \frac{(-1)^{(\ell-1)/2} (\ell-2)!!}{2^{(\ell+1)/2} ((\ell+1)/2)!}$ para $\ell$ impar:

$$A_\ell = \frac{3(2\ell+1)}{2 \alpha^\ell} \left[ \frac{(-1)^{(\ell-1)/2} (\ell-2)!!}{2^{(\ell+1)/2} ((\ell+1)/2)!} \right]$$

Solución final

La temperatura en el interior de la esfera es:

$$u(r, \theta) = \frac{1}{2} + \sum_{\ell=1, 3, 5, \dots}^{\infty} \left( \frac{3(2\ell+1)}{2 \alpha^\ell} \int_{0}^{1} P_\ell(x) \, dx \right) r^\ell P_\ell(\cos\theta)$$

\section{Problema 2}

Una esfera maciza posee un radio $R = 1$ a una temperatura $T(\theta, \phi) = \sin(\theta)e^{i\phi}$ en su superficie. La temperatura dentro de la esfera está dada por la ecuación de Laplace $\nabla^2 T(r, \theta, \phi) = 0$ y no existen fuentes caloríficas. Aplique separación de variables y halle $T(r, \theta, \phi)$ dentro de la esfera.

\subsection{Solución}

Para la ecuación de Laplace en coordenadas esféricas, cuando no hay dependencia radial externa, la solución general para el interior de una esfera (que debe ser finita en el origen) es:

$$T(r, \theta, \phi) = \sum_{\ell=0}^{\infty} \sum_{m=-\ell}^{\ell} A_{\ell m} r^\ell Y_{\ell m}(\theta, \phi)$$

Donde $Y_{\ell m}(\theta, \phi)$ son los armónicos esféricos.

Aplicación de la condición de contorno

En la superficie, donde $r = R = 1$, la solución debe coincidir con la temperatura dada:

$$T(1, \theta, \phi) = \sum_{\ell=0}^{\infty} \sum_{m=-\ell}^{\ell} A_{\ell m} (1)^\ell Y_{\ell m}(\theta, \phi) = \sin(\theta)e^{i\phi}$$

Identificación de modos

Sabemos que la función $\sin(\theta)e^{i\phi}$ es proporcional a uno de los armónicos esféricos. Específicamente, el armónico esférico $Y_{1,1}(\theta, \phi)$ está definido como:

$$Y_{1,1}(\theta, \phi) = -\sqrt{\frac{3}{8\pi}} \sin(\theta)e^{i\phi}$$

Por lo tanto, podemos expresar la condición de contorno como:

$$\sin(\theta)e^{i\phi} = -\sqrt{\frac{8\pi}{3}} Y_{1,1}(\theta, \phi)$$

Determinación de los coeficientes

Igualando los términos de la serie con la condición de contorno:

$$\sum_{\ell, m} A_{\ell m} Y_{\ell m}(\theta, \phi) = -\sqrt{\frac{8\pi}{3}} Y_{1,1}(\theta, \phi)$$

Por la ortogonalidad de los armónicos esféricos, el único coeficiente distinto de cero es $A_{1,1}$:

$$A_{1,1} = -\sqrt{\frac{8\pi}{3}}$$

Todos los demás $A_{\ell m} = 0$.

Resultado final

Sustituyendo el coeficiente $A_{1,1}$ y el radio $r$ en la solución general:

$$T(r, \theta, \phi) = A_{1,1} r^1 Y_{1,1}(\theta, \phi)$$

$$T(r, \theta, \phi) = \left( -\sqrt{\frac{8\pi}{3}} \right) r \left( -\sqrt{\frac{3}{8\pi}} \sin(\theta)e^{i\phi} \right)$$

Simplificando los términos, obtenemos la temperatura en el interior de la esfera:

$$T(r, \theta, \phi) = r \sin(\theta)e^{i\phi}$$




\bibliographystyle{plainurl}
\bibliography{bibliografia} 
\end{document}
